\documentclass[11pt]{article}
\usepackage[a4paper,margin=27mm]{geometry}
\usepackage{amsmath,amssymb,mathtools}
\usepackage{fontspec,unicode-math}
\setmainfont{Latin Modern Roman}
\setmathfont{Latin Modern Math}
\usepackage{xurl,hyperref}
\hypersetup{colorlinks=true,urlcolor=blue,linkcolor=blue}
\newcommand{\sref}[2]{\hyperlink{source:#1:#2}{[S#2]}}
\newcommand{\cataloguescope}[1]{\paragraph{Mathematical question.} #1}
\begin{document}
\section{Exceptional parameters without D(n)-quadruples}\label{statement}

\subsection{Definitions and mathematical statement}

\paragraph{Shared notation.}$\mathbb N=\{1,2,\ldots\}$, and $\mathbb Z,\mathbb Q,\mathbb R,\mathbb C$ denote the integers, rationals, reals and complex numbers. Any other conventions are specified locally.


For $n\in\mathbb Z$, a $D(n)$-$m$-tuple is a set $A=\{a_1,\ldots,a_m\}\subset\mathbb N$ of distinct positive integers such that $a_i a_j+n$ is the square of an integer for every $i<j$.
\paragraph{Statement recorded in the supplied source.}
Is it true that no $D(n)$-quadruple exists for each $n\in\{-3,3,5,8,12,20\}$? \sref{statement}{2}
\subsection{Short English statement}
For each of these six exceptional parameters, rule out or exhibit a set of four positive integers having the required pairwise square property.
\subsection{Sources}
\begin{description}
\item[\hypertarget{source:statement:1}{S1}] ULAM.AI and the UnsolvedMath Contributors, UnsolvedMath: A Curated Collection of Open Mathematics Problems, catalogue entry. CC BY 4.0. \url{https://huggingface.co/datasets/ulamai/UnsolvedMath}.
\item[\hypertarget{source:statement:2}{S2}] Andrej Dujella, Diophantine m-tuples, author-maintained survey (snapshot accessed 9 October 2026), Section 3.1, Theorem 3.2 and the exceptional set S. \url{https://raw.githubusercontent.com/dujella/dujella.github.io/master/dn.html}.
\end{description}
\label{end:statement}
\end{document}
